\hypertarget{appendix-b677a85ba76f}{}
\section*{Multiscale texture filtering}
\section{Multiscale texture filtering}\label{multiscale-texture-filtering}

\textbf{Curriculum depth:} Introduced · \textbf{First seen:} XR01 · \textbf{Latest development:} XR01

\subsection{Definition and mechanism}\label{definition-and-mechanism}

Mipmaps prefilter texture scales; anisotropy reflects unequal footprint directions. Gaussian pyramids store coarse structure; Laplacian levels store differences between scales. A texel is not an output pixel, and interpolation cannot ensure lost information is recovered.

\subsection{Prerequisites}\label{prerequisites}

\begin{itemize}
\tightlist
\item
  \href{}{Convolution and filter footprints (03-Visual-Computing-I, appendix-d5b3c692d9f0)} --- A weighting footprint establishes the filtered response, boundary effects and frequency transfer.
\item
  \hyperref[appendix-6a7869dd79a6]{Sampling aliasing and reconstruction} --- Allocation/filtering must respect what reduced samples can distinguish.
\end{itemize}

\subsection{Used by}\label{used-by}

\begin{itemize}
\tightlist
\item
  \hyperref[appendix-4192b2bdd21e]{Antialiasing and sample allocation}
\end{itemize}

\subsection{Recall and next depth}\label{recall-and-next-depth}

Explain a grazing-angle footprint and level log 2 rho. Pending: derivative-based GPU footprints and calibrated reconstruction. Written coverage is not a claim of learner mastery or accepted numbered-curriculum progress.

\subsection{Source and navigation}\label{source-and-navigation}

\href{Continuous\%20Notes/University\%20-\%20Worlds,\%20Images\%20and\%20Inference\%20-\%20Graphics,\%20Vision\%20and\%20Photography\#Multiresolution\%20work\%20must\%20preserve\%20the\%20right\%20information}{Expanded teaching explanation} · \href{}{Independent coverage (01-Foundations-I, appendix-fd090d982a15 § XR01 - Independent seeing and representation source)} · \href{}{Dependency Map (02-Foundations-II, appendix-4ad47bdd269a)}
